\sct{Non-Cartesian Coordinate Systems}
\ssc{General}
\sssc{Scale factors}
For a coordinate system, the scale factor $h_a$ of each coordinate $a$ is defined as
\[h_a=\norm{\pdv{\mb{r}}{a}}\]
And the scale factors are denoted with a vector of the same dimension as the space where each component is the scaling factor of that coordinate.
\sssc{Local basis}
For a coordinate system, the local basis is a position-dependent basis where each unit vector $\hat{\mb{e}}_a$ of each coordinate $a$ is defined as
\[\hat{\mb{e}}_a=\frac{1}{h_a}\pdv{\mb{r}}{a}\]
\sssc{Orthogonal coordinate system}
An orthogonal coordinate system is a coordinate system whose local basis is orthogonal.
\sssc{Line element}
For a coordinate system $(a_1,a_2,\ldots,a_n)$,
\[\dd{\mb{r}}=\sum_{i=1}^nh_{a_i}\hat{\mb{e}}_{a_i}\dd{a_i}.\]
\sssc{Lebesgue measure element}
For an orthogonal curvilinear coordinate system $(a_1,a_2,\ldots,a_n)$,
\[\dd{\mu}=\prod_{i=1}^nh_{a_i}\hat{\mb{e}}\dd{a_i}.\]
\sssc{Hypersurface elements}
For an orthogonal curvilinear coordinate system $(a_1,a_2,\ldots,a_n)$, the hypersurface element when $a_i$ is constants is
\[\dd{\mb{S}}=\pm\hat{\mb{e}}_{a_i}\prod_{\substack{j=1\\j\neq i}}^nh_{a_j}\prod_{\substack{j=1\\j\neq i}}^n\dd{a_j},\]
where $\pm$ is $+$ for right-handed and $-$ for left-handed.
\ssc{Polar coordinate system (極座標系)}
\sssc{Definition}
A polar coordinate system is a two-dimensional coordinate system that specifies point positions by a point called pole and a polar axis (a reference ray) that starts from the pole. The two polar coordinates $(r,\theta)$ are:
\bit
\item the radial distance $r\in\bbR_{\geq 0}$ defined as the distance of the point from the origin; and
\item the polar angle, azimuth, or azimuthal angle $\theta$ defined as any signed angle of the ray from the pole to the point from the polar axis with counterclockwise being positive and clockwise being negative.
\eit
For $r<0$, $(r,\theta)$ is defined as $(-r,\theta+\pi)$, which is not canonical.
\sssc{Conversion from and to Cartesian coordinate system}
If the pole and polar axis of a polar coordinate system is the origin and the positive direction of the $x$ axis of a Cartesian coordinate system, then:

To:
\[x=r\cos\theta,\]
\[y=r\sin\theta;\]
\[r=\sqrt{x^2+y^2},\]
From:
\[\theta=\operatorname{arctan2}(y,x)+2k\pi,\quad k\in\bbZ.\]
\sssc{Polar rectangle}
A polar rectangle is a region in $\bbR^2$ of the form
\[R=\{(r,\theta)\mid a\le r\le b\land\alpha\le\theta\le\beta\}\]
where $0\le a\le b$ and $0\le\beta-\alpha\le 2\pi$.
\sssc{Local orthogonal basis}
Radial unit vector $\hat{\mb{e}}_r$ is $(\cos\theta,\sin\theta)$ (Cartesian). Angular (transverse) unit vector $\hat{\mb{e}}_{\theta}$ is $(-\sin\theta,\cos\theta)$ (Cartesian).
\sssc{Scaling factors}
\[(h_r,h_{\theta})=(1,r)\]
\sssc{Line element}
\[\dd{\mb{r}}=\dd{r}\hat{\mb{e}}_r+r\dd{\theta}\hat{\mb{e}}_{\theta}\]
\bpr
\[\ba
\dd{\mb{r}}&=\dd{(r\hat{\mb{e}}_r)}\\
&=\dd{r}\hat{\mb{e}}_r+r\dd{\hat{\mb{e}}_r}\\
&=\dd{r}\hat{\mb{e}}_r+r\dd{\theta}\hat{\mb{e}}_{\theta}
\ea\]
\epr
\sssc{Area element}
\[\dd{A}=r\dd{r}\dd{\theta}\]
\bpr
\[T\colon(r,\theta)\mapsto(x,y)=(r\cos\theta,r\sin\theta).\]
\[J_T=\begin{pmatrix}
\cos\theta & -r\sin\theta\\
\sin\theta & r \cos\theta
\end{pmatrix}\]
\[\abs{\det(J_T)}=r\]
\epr
\tb{Riemann integral}:

Let
\[R=\{(r,\theta)\mid a\le r\le b\land\alpha\le\theta\le\beta\},\quad0\le a\le b\land0\le\beta-\alpha\le 2\pi.\]
Then,
\[\iint_Rf(x,y)\dd{A}=\int_{\alpha}^{\beta}\int_a^bf(r\cos\theta,r\sin\theta)r\dd{r}\dd{\theta}.\]
\bpr
For each
\[R_{ij}=[r_{i-1},r_i]\times[\theta_{i-1},\theta_i],\]
let
\[r_i^*=\frac{1}{2}(r_{i-1}+r_i)\]
\[\theta_i^*=\frac{1}{2}(\theta_{i-1}+\theta_i)\]
\[\Delta A_i=r_i^*\Delta r\Delta\theta\]
\[\sum_{i=1}^m\sum_{j=1}^nf(r_i^*\cos\theta_j^*,r_i^*\sin\theta_j^*)\Delta A_i=\sum_{i=1}^m\sum_{j=1}^nf(r_i^*\cos\theta_j^*,r_i^*\sin\theta_j^*)r_i^*\Delta r\Delta\theta\]
\epr
\ssc{ISO 31-11 cylindrical coordinate system (圓柱座標系)}
\sssc{Definition}
A cylindrical coordinate system is a three-dimensional coordinate system that specifies point positions by a main axis (a chosen directed line) and an auxiliary axis (a reference ray) with the intersection of the two axes called the origin and the plane passing through the origin and perpendicular to the main axis equipped with polar coordinate system with the polar axis being the auxiliary axis called the reference plane. The three cylindrical coordinates $(r,\theta,z)$ are:
\bit
\item the radial distance $r\in\bbR_{\geq 0}$ defined as the perpendicular distance of the point from the main axis;
\item the azimuth or azimuthal angle $\theta\in(-\pi,\pi]$ defined as the polar angle of the projection of the point on the reference plane and undefined on the main axis; and
\item the axial coordinate or height $z\in\bbR_{\geq 0}$ defined as the signed distance of the point from the reference plane where the sign is defined to be positive if the direction is the same as the main axis and negative if the direction is the negation of the direction of the main axis.
\eit
\sssc{Conversion from and to Cartesian coordinate system}
If the main axis and the auxiliary axis of a cyclindrical coordinate system is the $z$ axis and the positive direction of the $x$ axis of a Cartesian coordinate system, then:

To:
\[x=r\cos\theta,\]
\[y=r\sin\theta,\]
\[z=z;\]
From:
\[r=\sqrt{x^2+y^2},\]
\[\theta=\begin{cases}\sgn(y)\arccos\frac{x}{\sqrt{x^2+y^2}},\quad&y\neq 0\lor x>0\\\pi,\quad&y=0\land x<0\end{cases},\]
\[z=z.\]
\sssc{Local orthogonal basis}
Radial unit vector $\hat{\mb{e}}_r$ is $(\cos\theta,\sin\theta,0)$ (Cartesian). Azimuthal unit vector $\hat{\mb{e}}_{\theta}$ is $(-\sin\theta,\cos\theta,0)$ (Cartesian). Axial unit vector $\hat{\mb{e}}_z$ is $(0,0,1)$ (Cartesian).
\sssc{Scaling factors}
\[(h_r,h_{\theta},h_z)=(1,r,1)\]
\sssc{Line element}
\[\dd{\mb{r}}=\dd{r}\hat{\mb{e}}_r+r\dd{\theta}\hat{\mb{e}}_{\theta}+\dd{z}\hat{\mb{e}}_z\]
\sssc{Volume element}
\[\dd{V}=r\dd{r}\dd{\theta}\dd{z}\]
\bpr
\[T\colon(r,\theta,z)\mapsto(x,y,z)=(r\cos\theta,r\sin\theta,z).\]
\[J_T=\begin{pmatrix}
\cos\theta & -r\sin\theta & 0\\
\sin\theta & r \cos\theta & 0\\
0 & 0 & 1
\end{pmatrix}\]
\[\abs{\det(J_T)}=r\]
\epr
\sssc{Surface element}
$r=$ constant:
\[\dd{\mb{S}}=\hat{\mb{e}}_rr\dd{\theta}\dd{z}\]
$\theta=$ constant:
\[\dd{\mb{S}}=\hat{\mb{e}}_{\theta}\dd{r}\dd{z}\]
$z=$ constant:
\[\dd{\mb{S}}=\hat{\mb{e}}_zr\dd{r}\dd{\theta}\]
\ssc{ISO 80000-2, inclination-azimuth, or physics convention spherical coordinate system (球座標系)}
A spherical coordinate system is a three-dimensional coordinate system that specifies point positions by a polar axis (a chosen directed line) and an auxiliary axis (a reference ray) with the intersection of the two axes called the origin and the plane passing through the origin and perpendicular to the polar axis equipped with polar coordinate system with the polar axis being the auxiliary axis called the reference plane. The three spherical coordinates $(r,\theta,\varphi)$ are:
\bit
\item the radial distance or radius $r\in\bbR_{\geq 0}$ defined as the distance of the point from the origin;
\item the polar angle, inclination, or zenith angle $\theta\in[0,\pi]$ defined as the angle between the ray from the origin to the point and the positive half of the polar axis and undefined at the origin; and
\item the azimuth or azimuthal angle $\varphi\in(-\pi,\pi]$ defined as the polar angle of the projection of the point on the reference plane and undefined on the polar axis (of the coordinate system instead of the reference plane).
\eit
\ti{Calculus: Early Transcendentals} by James Stewart exchanges the definition of $\theta$ and $\varphi$ and uses $\rho$ to denote $r$.
\sssc{Conversion from and to Cartesian coordinate system}
If the polar axis and the auxiliary axis of a spherical coordinate system is the $z$ axis and the positive direction of the $x$ axis of a Cartesian coordinate system, then:

To:
\[x=r\sin\theta\cos\phi,\]
\[y=r\sin\theta\sin\phi,\]
\[z=r\cos\theta;\]
From:
\[r=\sqrt{x^2+y^2+z^2},\]
\[\theta=\arccos\frac{z}{r},\]
\[\varphi=\begin{cases}\sgn(y)\arccos\frac{x}{\sqrt{x^2+y^2}},\quad&y\neq 0\lor x>0\\\pi,\quad&y=0\land x<0\end{cases}.\]
\sssc{Conversion from and to cylindrical coordinate system}
If the polar axis and the auxiliary axis of a spherical coordinate system is the main axis and the auxiliary axis of a cyclindrical coordinate system, then:

To:
\[r=r\sin\theta,\]
\[\theta=\varphi,\]
\[z=r\cos\theta;\]
From:
\[r=\sqrt{r^2+z^2},\]
\[\theta=\arctan\frac{r}{z},\]
\[\varphi=\theta.\]
\sssc{Local orthogonal basis}
Radial unit vector $\hat{\mb{e}}_r$ is $(\sin\theta\cos\varphi,\sin\theta\sin\varphi,\cos\theta)$ (Cartesian). Polar unit vector $\hat{\mb{e}}_{\theta}$ is $(\cos\theta\cos\varphi,\cos\theta\sin\varphi,-\sin\theta)$ (Cartesian). Azimuthal unit vector $\hat{\mb{e}}_{\varphi}$ is $(-\sin\varphi,\cos\varphi,0)$ (Cartesian).
\sssc{Scale factor}
\[(h_r,h_{\theta},h_{\varphi})=(1,r,r\sin\theta)\]
\sssc{Line element}
\[\dd{\mb{r}}=\dd{r}\hat{\mb{e}}_r+r\dd{\theta}\hat{\mb{e}}_{\theta}+r\sin\theta\dd{\varphi}\hat{\mb{e}}_{\varphi}\]
\sssc{Volume element}
\[\dd{V}=r^2\sin\theta\dd{r}\dd{\theta}\dd{\varphi}\]
\sssc{Surface elements}
$r=$ constant:
\[\dd{\mb{S}}=\hat{\mb{e}}_rr^2\sin\theta\dd{\theta}\dd{\varphi}\]
$\theta=$ constant:
\[\dd{\mb{S}}=\hat{\mb{e}}_{\theta}r\sin\theta\dd{r}\dd{\varphi}\]
$\varphi=$ constant:
\[\dd{\mb{S}}=\hat{\mb{e}}_{\varphi}r\dd{r}\dd{\theta}\]
